\scanpage{4}
% Reused typescript is upside down in the supplied scan; text order follows its reading orientation.
\mspage{3}
\begin{center}-3.2-\end{center}
\margin{\unsure{11}}
\[
\varpi:\operatorname{SL}(n+1)\longrightarrow\operatorname{PGL}(n)
\]
is the canonical isogeny, then with respect to the induced action of
$\operatorname{SL}(n+1)$, $o_{\mathbb P}(1)$ admits an $\operatorname{SL}(n+1)$-linearization. This follows
because $\operatorname{SL}(n+1)$ acts on the affine cone over $\mathbb P_n$, compatibly with
the given action of $\operatorname{PGL}(n)$ on $\mathbb P_n$. This linearization induces an
$\operatorname{SL}(n+1)$-linearization of $L_i$ for each $i$ (for the induced action of
$\operatorname{SL}(n+1)$), and therefore such a linearization for each product
\[
\bigotimes_{i=0}^{m}L_i^{\otimes r_i}.
\]
On the other hand, the $(n+1)^{\mathrm{st}}$ power of the $\operatorname{SL}(n+1)$-linearization of
$\underline{o}_{\mathbb P}(1)$, which is a linearization of $\underline{o}_{\mathbb P}(n+1)$, is induced, via $\varpi$,
from a $\operatorname{PGL}(n)$-linearization (cf. Chapter 1, §2). Therefore, the
$\operatorname{SL}(n+1)$-linearization of
\[
\bigotimes_{i=0}^{m}L_i^{r_i}
\]
is a $\operatorname{PGL}(n)$-linearization if $n+1$ divides each $r_i$.

In any case, by means of these linearizations we can define in-
variant sections of all the sheaves $\bigotimes_{i=0}^{m}L_i^{r_i}$. To construct such
invariant sections, let $X_0,\ldots,X_n$ be the canonical sections of $\underline{o}_{\mathbb P}(1)$
on $\mathbb P_n$. Let
\[
X_i^{(j)}=p_j^*(X_i)
\]
be the induced sections of $L_j$.

\uline{Definition} \uline{3.1}: For all sequences $\alpha_0,\ldots,\alpha_n$ of integers such that
